\section{Introduction}\label{sec:intro}

How much of a child's economic position is inherited is usually summarised by one number, the intergenerational
elasticity of income (IGE) or its rank analogue, measured when the children are in their early thirties. This paper
shows that the number depends heavily on the age at which children are observed, and that it depends on age
differently for Black and White Americans. In the Panel Study of Income Dynamics (PSID), the elasticity of sons' own
earnings with respect to their parents' income is 0.28 at ages 25--29 and 0.49 at 40--44, and the rise is concentrated
among White sons. The White--Black gap in the elasticity is zero at 25--29 ($-0.04$), 0.10 in the 31--37 window used by
the benchmark administrative estimates of \citet{chetty2020race}, and 0.29 at 40--49 (Figure~\ref{fig:headline}). Measured
by potential experience rather than age, the gap is zero throughout the first ten years of the career and opens
after.\footnote{Child-year regressions within five-year age bands, Table~\ref{tab:ige_sons}; 31--37: 0.099$^{**}$
(0.048); 40--44 and 45--49: 0.294$^{***}$ (0.067) and 0.299$^{***}$ (0.088). By potential experience: $-0.121$ at 0--4
years, 0.043 at 5--9, 0.220$^{***}$ at 15--19 and 0.272$^{***}$ at 20--24 (Appendix Table~\ref{tab:experience}). Code:
\texttt{IGE/26.IGEByAge.py} and \texttt{IGE/46.AgeProfileRound2.py}.} A window that ends at 37 sees a third of the racial
gap in persistence that is present in midlife; a window in the late twenties sees none of it.

The racial half of this fact is estimated on 500 to 1,700 Black sons per age band in the PSID, and I state it at its
precision. With one year of parental income, the one measure four public data sets share, the gap opens by 0.09 log
points between the late twenties and the early forties (0.086$^{***}$, s.e.\ 0.024); corrected by each data set's own
reliability ratio of one year of parental income, by 0.19 (0.191$^{***}$, s.e.\ 0.046; Table~\ref{tab:threedata}). The
PSID estimate survives the standard objections: selective attrition, period effects, the dollar-weighted elasticity of
\citet{mitnik2020intergenerational}, and an experience clock in place of age all leave its opening between 0.20 and 0.26
(Table~\ref{tab:gaprobust}). In the fourth data set, the children of the NLSY79 mothers, the gap is zero at 25--34 as in
the others and does not open, but it is measured at 35--44 on fewer than 300 sons of each race, born to young mothers. The object of
the paper is sons' earnings and wages. Their family income, the object of \citet{chetty2020race}, shows the same profile
with a smaller and more fragile opening: it is not significant in a fixed calendar window (0.084, s.e.\ 0.084), and among sons
of parents below the median it is 0.023 (s.e.\ 0.086).

The log gap is the age timing of a known fact. Black sons of well-off parents are far more likely than White sons to
fall down the distribution \citep{hertz2005rags, bhattacharya2011nonparametric, mazumder2014black, chetty2020race}, and
racial differences have shaped American mobility over the century \citep{jacome2025mobility}. This paper shows when the
difference appears: it is absent in the late twenties and builds through the thirties as White sons' earnings spread out
with parental income. Ranks of family income normalise that spread away, and sons' rank--rank slopes of family income are
similar by race at every age, as \citet{chetty2020race} find at 31--37.

What opens is the transmission of hourly wages, not hours, shocks or marriage (Section~\ref{sec:contains}), and it is
built at 25--34. The job a son holds then accounts for half of the racial divergence (0.49), and access to the
well-paid jobs whose wages keep growing is unequal by parental income---20 percent of sons from the top third of
parental incomes start in one, 5 percent from the bottom third---and by race at equal parental income (4.2 percentage
points). The economic channel is the return to schooling. It rises with age much faster for White than for Black sons
(a difference of 0.058$^{***}$ in the rise of the earnings return), and the difference is the same within the job a son
started in, net of the calendar rise of the college premium, and in the NLSY79 holding measured ability fixed. Education
accounts for 0.44 of the racial divergence in earnings persistence.\footnote{Section~\ref{sec:careers},
Tables~\ref{tab:earlyjob}, \ref{tab:access}, \ref{tab:schooling} and \ref{tab:withincell}. Within the early job cell
the difference in the rise is 0.061$^{***}$ (0.018); with schooling $\times$ calendar trend 0.063$^{**}$ (0.027); in
the NLSY79 with the AFQT 0.021$^{**}$ (0.009). Code: \texttt{IGE/35.EarlyIndustryLatePayoff.py},
\texttt{IGE/42.SchoolingAndJobGrowth.py}, \texttt{IGE/45.WithinCellIdentification.py}.} Both the schooling and the AFQT
returns rise with age, which is the complementarity of schooling with human capital accumulated on the job rather than
employers learning about ability \citep{altonji2001employer, lange2007speed}. Fathers' jobs add a separate channel: sons
of richer White parents converge on their fathers' occupations after 35, and convergence on the father's exact job is
the inheritance of a family business (0.065$^{***}$ for sons of self-employed fathers, zero for others).

The facts bear on how policies and shocks affect racial disparities (Section~\ref{sec:policy}). Long-run
evaluations of childhood policies---Food Stamps, school finance, the EITC, Moving to Opportunity---measure adult
earnings in the late twenties or early thirties \citep{hoynes2016long, jackson2016effects, bastian2018long,
chetty2016effects}, before the racial gap in persistence has opened, so their racial comparisons are scored too early.
Two shocks to the start of the career do not produce the gap: recessions at labor-market entry leave average access to
good jobs unchanged and the racial gap opens among sons who entered tight markets too, although they leave a longer
wage scar on Black sons; and public employment, where Black sons start more often, carries none of the
divergence.\footnote{Section~\ref{sec:policy}, Tables~\ref{tab:entry} and \ref{tab:openinggroups}. Shocks to the household a
son grows up in---his parents' job losses and the family safety net---are the subject of a companion paper; a shock to
the son's own career is this paper's.} These facts leave four explanations: at equal background, Black sons may have less
access to the jobs in which human capital accumulates fast, a lower pay-off of schooling in accumulation, a different
return to schooling at entry, or a lower pay-off of parental income in accumulation. Section~\ref{sec:model} nests all
four in one life-cycle model of careers---jobs that differ in accumulation, reached through offers and search
frictions, job loss that falls with age and parental income, and accumulation that depends on schooling and
background---and estimates every combination of the first three, with and without the fourth, on the same 66 moments. The data select one model. Parental
income raises the growth of human capital for White sons (0.019$^{***}$ a year per log point, $p = 0.009$) and not for
Black sons (difference $-0.019^{***}$, $p < 0.001$), a channel seen with a multi-year measure of parental income and
absorbed by schooling with a one-year measure, in the PSID as in the NLSY79 and NLSY97 (Section~\ref{sec:model}); for
Black sons schooling buys a higher wage at entry and slower
growth; and a racial difference in offers of good jobs is rejected: the lower Black share in good jobs at equal
parental income is reproduced without one. The model reproduces the rising White elasticity
and the flat Black one, and switching off the racial difference in the pay-off of parental income closes three fifths
of the gap at 45--49 and two fifths of its opening since 25--29. Equalising schooling, access across parental income or entry conditions closes less than a tenth:
the gap at 45 is built by what background and schooling buy in growth, not by where careers start. The paper also makes a falsifiable prediction for administrative data: for the cohorts born 1978--83, the racial
gap in the log elasticity of individual earnings at 43--47 should be 0.25--0.30, against about 0.10 at 31--37.

\paragraph{Contributions and literature.} The paper contributes to four literatures. First, the measurement of
intergenerational persistence: short snapshots of parental income understate persistence
\citep{solon1992intergenerational, mazumder2005fortunate, mazumder2016estimating}, current income early in the career
understates the dispersion of lifetime income \citep{haider2006life, grawe2006lifecycle}, heterogeneous profiles leave a
life-cycle bias that no single correction removes \citep{nybom2016heterogeneous, nybom2017biases}, ranks stabilise early
\citep{chetty2014where, chetty2014still}, and the elasticity can be redefined on expected income
\citep{mitnik2015new, mitnik2020intergenerational}. I show that the life-cycle bias differs by race. Second, race and
mobility \citep{hertz2005rags, bhattacharya2011nonparametric, mazumder2014black, chetty2020race, bayer2018divergent,
jacome2025mobility}: I date the racial difference in log persistence to the first decade of the career and reconcile it
with stable rank gaps. Third, careers and the Black--White wage gap: wages grow through job mobility and on-the-job
accumulation \citep{topel1992job, altonji2013modeling, bagger2014tenure}; the racial gap widens with experience
\citep{oettinger1996statistical, lang2012racial}; premarket skills, schooling and employer learning shape it
\citep{neal1996role, altonji2001employer, pinkston2006test, lange2007speed, arcidiacono2010beyond, lang2011education};
search frictions can differ by race \citep{bowlus2002discrimination}; promotion is slower for equally eligible Black
workers \citep{hegde2026racial}; education premiums rise over the life cycle \citep{bhuller2017life}; and lifetime
inequality is largely set at entry \citep{huggett2011sources}; a calibrated life-cycle model with learning by doing
attributes the widening of the racial gap to premarket differences in the ability to accumulate human capital, holding
accumulation race-neutral given the AFQT \citep{rauh2023black}. I connect these to family background and show that the
rise of the return to schooling with age is steeper for White men within the same job. Fourth, family networks and the
inheritance of jobs \citep{corak2011intergenerational, kramarz2014strong, staiger2022, laband1992self,
fagereng2021wealthy}: following the father pays at entry, not in growth, and late convergence on the father's exact job
is business inheritance. Structural models of intergenerational transmission concentrate family background before
entry \citep{lee2019intergenerational, caucutt2020early, abbott2019education, bolt2024intergenerational}; the evidence
here requires that background, and race, also shape what happens after entry.

Section~\ref{sec:data} describes the data, Section~\ref{sec:ageprofile} the age profile and its precision,
Section~\ref{sec:contains} what opens, Section~\ref{sec:careers} where it is built, Section~\ref{sec:fathers} fathers'
jobs, Section~\ref{sec:policy} shocks and policies, and Section~\ref{sec:model} the model and the policy experiments.
